\documentclass[11pt, letterpaper]{article}

\usepackage[margin=1in]{geometry}
\usepackage{amsmath, amssymb}
\usepackage{booktabs}
\usepackage{longtable}
\usepackage{array}
\usepackage{multirow}
\usepackage{xcolor}
\usepackage{colortbl}
\usepackage{hyperref}
\usepackage{titlesec}
\usepackage{parskip}
\usepackage{enumitem}
\usepackage{fancyhdr}
\usepackage{graphicx}
\usepackage{caption}

% Shaded zero cell shorthand (global so longtable can use it across pages)
\definecolor{zerogray}{RGB}{220, 220, 220}
\newcommand{\Z}{\cellcolor{zerogray}0}
\newcolumntype{C}{>{\centering\arraybackslash}p{0.9cm}}

% Section formatting
\titleformat{\section}{\large\bfseries}{}{0em}{}[\titlerule]
\titleformat{\subsection}{\normalsize\bfseries}{}{0em}{}
\titlespacing{\section}{0pt}{18pt}{8pt}
\titlespacing{\subsection}{0pt}{12pt}{4pt}

% Header / footer
\pagestyle{fancy}
\fancyhf{}
\rhead{\small\color{gray}IEOR E4004 --- Spring 2026 $\mid$ Project II}
\lhead{\small\color{gray}Scheduling the NBA}
\rfoot{\small\thepage}
\renewcommand{\headrulewidth}{0.4pt}

% Hyperlink style
\hypersetup{
  colorlinks=false,
  hidelinks,
}

% ─────────────────────────────────────────────────────────────────────────────
\begin{document}

% ── Title block ──────────────────────────────────────────────────────────────
\begin{center}
  {\LARGE\bfseries IEOR E4004 --- Project II}\\[6pt]
  {\Large\bfseries Scheduling the NBA}\\[10pt]
  {\normalsize Spring 2026}\\[6pt]
  {\normalsize
    Angel Huang (ah4422) \quad Cole Zemel (ccz2107) \quad
    Nathan Layani (ndl2127) \quad Suyi Li (sl5867)}
\end{center}

\vspace{6pt}
\noindent\rule{\linewidth}{1.5pt}
\vspace{4pt}

% ── Table of Contents ────────────────────────────────────────────────────────
\tableofcontents
\newpage

% =============================================================================
\section{Question 1 --- Schedule Analysis}
% =============================================================================

For each team $i$ we extracted the following four items from the preliminary
schedule:

\begin{itemize}[nosep]
  \item[(a)] All dates when team $i$ played at home.
  \item[(b)] For each opponent $j$, the number of times team $i$ hosted $j$ at
             home.
  \item[(c)] For each opponent $j$, the number of times team $i$ played away at
             $j$'s arena.
  \item[(d)] All dates when team $i$ played away.
\end{itemize}

% ── (a) & (d) ────────────────────────────────────────────────────────────────
\subsection{(a) \& (d) Home and Away Dates}

The table below lists each team's home game dates (a) and away game dates (d).
The schedule spans November~1 -- December~25, 2025 across 16 dates, with 8
games per date (128 total).

\begingroup
\small
\setlength{\tabcolsep}{4pt}
\begin{longtable}{p{3.2cm} p{6.8cm} p{5.2cm}}
\toprule
\textbf{Team} & \textbf{(a) Home Dates} & \textbf{(d) Away Dates} \\
\midrule
\endhead
\bottomrule
\endfoot
Atlanta Hawks       & Nov 03, 07, 15, 17, 19, 23, 27, 28, 29; Dec 25 & Nov 01, 05, 11, 13, 21; Dec 01 \\[2pt]
Boston Celtics      & Nov 01, 07, 13, 19, 23, 28 & Nov 03, 05, 11, 15, 17, 21, 27, 29; Dec 01, 25 \\[2pt]
Brooklyn Nets       & Nov 01, 05, 11, 13, 21, 29; Dec 01 & Nov 03, 07, 15, 17, 19, 23, 27, 28; Dec 25 \\[2pt]
Chicago Bulls       & Nov 01, 03, 13, 21, 23, 27, 28, 29; Dec 01 & Nov 05, 07, 11, 15, 17, 19; Dec 25 \\[2pt]
Cleveland Cavaliers & Nov 03, 07, 17, 19, 21, 23, 28, 29; Dec 01, 25 & Nov 01, 05, 11, 13, 15, 27 \\[2pt]
Dallas Mavericks    & Nov 03, 07, 15, 17, 19, 23, 27, 28; Dec 25 & Nov 01, 05, 11, 13, 21, 29; Dec 01 \\[2pt]
Denver Nuggets      & Nov 05, 07, 11, 15, 17, 19, 21; Dec 25 & Nov 01, 03, 13, 23, 27, 28, 29; Dec 01 \\[2pt]
Golden State Warriors & Nov 03, 05, 11, 15, 17, 27; Dec 25 & Nov 01, 07, 13, 19, 21, 23, 28, 29; Dec 01 \\[2pt]
Houston Rockets     & Nov 03, 11, 13, 15, 19, 21, 23, 27; Dec 25 & Nov 01, 05, 07, 17, 28, 29; Dec 01 \\[2pt]
Los Angeles Lakers  & Nov 03, 05, 11, 15, 17, 19; Dec 25 & Nov 01, 07, 13, 21, 23, 27, 28, 29; Dec 01 \\[2pt]
Miami Heat          & Nov 01, 05, 13, 17, 27, 28, 29; Dec 01 & Nov 03, 07, 11, 15, 19, 21, 23; Dec 25 \\[2pt]
Milwaukee Bucks     & Nov 01, 07, 13, 21, 23, 27, 28, 29; Dec 01 & Nov 03, 05, 11, 15, 17, 19; Dec 25 \\[2pt]
New York Knicks     & Nov 01, 05, 11, 13, 15; Dec 01 & Nov 03, 07, 17, 19, 21, 23, 27, 28, 29; Dec 25 \\[2pt]
Philadelphia 76ers  & Nov 01, 05, 07, 11, 21, 29; Dec 01 & Nov 03, 13, 15, 17, 19, 23, 27, 28; Dec 25 \\[2pt]
Phoenix Suns        & Nov 03, 11, 13, 15, 19, 21, 23; Dec 25 & Nov 01, 05, 07, 17, 27, 28, 29; Dec 01 \\[2pt]
Toronto Raptors     & Nov 01, 05, 07, 17, 27, 28, 29; Dec 01 & Nov 03, 11, 13, 15, 19, 21, 23; Dec 25 \\
\end{longtable}
\endgroup

% ── (b) Home matchup counts ──────────────────────────────────────────────────
\subsection{(b) Home Matchup Counts --- $c^H(i,j)$}

Each cell shows the number of times the row team $i$ hosted the column team $j$
at home. A value of 0 (shaded) indicates no home game between that pair;
``---'' marks a team against itself.

\begingroup
\footnotesize
\setlength{\tabcolsep}{3pt}

\begin{longtable}{p{3.0cm}CCCCCCCCCCCCCCCC}
\toprule
\textbf{Team $\backslash$ Opp.} & \textbf{ATL} & \textbf{BOS} & \textbf{BKN} & \textbf{CHI} & \textbf{CLE} & \textbf{DAL} & \textbf{DEN} & \textbf{GSW} & \textbf{HOU} & \textbf{LAL} & \textbf{MIA} & \textbf{MIL} & \textbf{NYK} & \textbf{PHI} & \textbf{PHX} & \textbf{TOR} \\
\midrule
\endhead
\bottomrule
\endfoot
Atlanta Hawks        & ---  & 1    & \Z   & 1    & \Z   & 1    & \Z   & 1    & \Z   & 1    & 1    & 1    & 1    & \Z   & 1    & 1    \\
Boston Celtics       & \Z   & ---  & \Z   & \Z   & \Z   & 1    & \Z   & 1    & \Z   & \Z   & 1    & 1    & \Z   & 1    & 1    & \Z   \\
Brooklyn Nets        & 1    & 1    & ---  & \Z   & \Z   & 1    & 1    & 1    & \Z   & \Z   & 1    & \Z   & \Z   & 1    & \Z   & \Z   \\
Chicago Bulls        & 1    & 1    & 1    & ---  & \Z   & \Z   & 1    & 1    & 1    & 1    & \Z   & 1    & \Z   & \Z   & \Z   & 1    \\
Cleveland Cavaliers  & 1    & 1    & 1    & 1    & ---  & \Z   & 1    & \Z   & 1    & \Z   & 1    & 1    & 1    & \Z   & \Z   & 1    \\
Dallas Mavericks     & \Z   & \Z   & \Z   & 1    & 1    & ---  & 1    & 1    & 1    & 1    & 1    & \Z   & \Z   & 1    & \Z   & 1    \\
Denver Nuggets       & 1    & 1    & 1    & \Z   & \Z   & \Z   & ---  & \Z   & 1    & 1    & 1    & \Z   & 1    & 1    & \Z   & \Z   \\
Golden State Warriors & \Z  & 1    & \Z   & \Z   & 1    & \Z   & 1    & ---  & 1    & \Z   & \Z   & 1    & 1    & 1    & \Z   & \Z   \\
Houston Rockets      & 1    & 1    & 1    & \Z   & \Z   & \Z   & \Z   & \Z   & ---  & 1    & \Z   & 1    & 1    & 1    & 1    & 1    \\
Los Angeles Lakers   & \Z   & 1    & 1    & \Z   & 1    & \Z   & \Z   & 1    & \Z   & ---  & \Z   & \Z   & 1    & \Z   & 1    & 1    \\
Miami Heat           & \Z   & \Z   & \Z   & 1    & \Z   & 1    & \Z   & 1    & 1    & 1    & ---  & \Z   & 1    & 1    & 1    & \Z   \\
Milwaukee Bucks      & \Z   & \Z   & 1    & \Z   & 1    & 1    & 1    & \Z   & \Z   & 1    & 1    & ---  & 1    & 1    & 1    & \Z   \\
New York Knicks      & \Z   & 1    & 1    & 1    & \Z   & 1    & \Z   & \Z   & \Z   & 1    & \Z   & \Z   & ---  & \Z   & \Z   & 1    \\
Philadelphia 76ers   & 1    & \Z   & \Z   & 1    & 1    & \Z   & \Z   & \Z   & \Z   & 1    & \Z   & \Z   & 1    & ---  & 1    & 1    \\
Phoenix Suns         & \Z   & \Z   & 1    & 1    & 1    & 1    & 1    & 1    & \Z   & \Z   & \Z   & \Z   & 1    & 1    & ---  & \Z   \\
Toronto Raptors      & \Z   & 1    & 1    & \Z   & \Z   & \Z   & 1    & 1    & 1    & \Z   & 1    & 1    & \Z   & \Z   & 1    & ---  \\
\end{longtable}
\endgroup

% ── (c) Away matchup counts ──────────────────────────────────────────────────
\subsection{(c) Away Matchup Counts --- $c^A(i,j)$}

Each cell shows the number of times the row team $i$ played away at the column
team $j$'s arena. Values of 0 are shaded.

\begingroup
\footnotesize
\setlength{\tabcolsep}{3pt}

\begin{longtable}{p{3.0cm}CCCCCCCCCCCCCCCC}
\toprule
\textbf{Team $\backslash$ Opp.} & \textbf{ATL} & \textbf{BOS} & \textbf{BKN} & \textbf{CHI} & \textbf{CLE} & \textbf{DAL} & \textbf{DEN} & \textbf{GSW} & \textbf{HOU} & \textbf{LAL} & \textbf{MIA} & \textbf{MIL} & \textbf{NYK} & \textbf{PHI} & \textbf{PHX} & \textbf{TOR} \\
\midrule
\endhead
\bottomrule
\endfoot
Atlanta Hawks        & ---  & \Z   & 1    & 1    & 1    & \Z   & 1    & \Z   & 1    & \Z   & \Z   & \Z   & \Z   & 1    & \Z   & \Z   \\
Boston Celtics       & 1    & ---  & 1    & 1    & 1    & \Z   & 1    & 1    & 1    & 1    & \Z   & \Z   & 1    & \Z   & \Z   & 1    \\
Brooklyn Nets        & \Z   & \Z   & ---  & 1    & 1    & \Z   & 1    & \Z   & 1    & 1    & \Z   & 1    & 1    & \Z   & 1    & 1    \\
Chicago Bulls        & 1    & \Z   & \Z   & ---  & 1    & 1    & \Z   & \Z   & \Z   & \Z   & 1    & \Z   & 1    & 1    & 1    & \Z   \\
Cleveland Cavaliers  & \Z   & \Z   & \Z   & \Z   & ---  & 1    & \Z   & 1    & \Z   & 1    & \Z   & 1    & \Z   & 1    & 1    & \Z   \\
Dallas Mavericks     & 1    & 1    & 1    & \Z   & \Z   & ---  & \Z   & \Z   & \Z   & \Z   & 1    & 1    & 1    & \Z   & 1    & \Z   \\
Denver Nuggets       & \Z   & \Z   & 1    & 1    & 1    & 1    & ---  & 1    & \Z   & \Z   & \Z   & 1    & \Z   & \Z   & 1    & 1    \\
Golden State Warriors & 1   & 1    & 1    & 1    & \Z   & 1    & \Z   & ---  & \Z   & 1    & 1    & \Z   & \Z   & \Z   & 1    & 1    \\
Houston Rockets      & \Z   & \Z   & \Z   & 1    & 1    & 1    & 1    & 1    & ---  & \Z   & 1    & \Z   & \Z   & \Z   & \Z   & 1    \\
Los Angeles Lakers   & 1    & \Z   & \Z   & 1    & \Z   & 1    & 1    & \Z   & 1    & ---  & 1    & 1    & 1    & 1    & \Z   & \Z   \\
Miami Heat           & 1    & 1    & 1    & \Z   & 1    & 1    & 1    & \Z   & \Z   & \Z   & ---  & 1    & \Z   & \Z   & \Z   & 1    \\
Milwaukee Bucks      & 1    & 1    & \Z   & 1    & 1    & \Z   & \Z   & 1    & 1    & \Z   & \Z   & ---  & \Z   & \Z   & \Z   & 1    \\
New York Knicks      & 1    & \Z   & \Z   & \Z   & 1    & \Z   & 1    & 1    & 1    & 1    & 1    & 1    & ---  & 1    & 1    & \Z   \\
Philadelphia 76ers   & \Z   & 1    & 1    & \Z   & \Z   & 1    & 1    & 1    & 1    & \Z   & 1    & 1    & \Z   & ---  & 1    & \Z   \\
Phoenix Suns         & 1    & 1    & \Z   & \Z   & \Z   & \Z   & \Z   & \Z   & 1    & 1    & 1    & 1    & \Z   & 1    & ---  & 1    \\
Toronto Raptors      & 1    & \Z   & \Z   & 1    & 1    & 1    & \Z   & \Z   & 1    & 1    & \Z   & \Z   & 1    & 1    & \Z   & ---  \\
\end{longtable}
\endgroup

% =============================================================================
\section{Question 2 --- Integer Programming Formulation}
% =============================================================================

We formulate a feasibility Integer Program (IP) whose feasible solutions are
exactly the valid schedules --- those that preserve all structural properties of
the original draft while potentially reassigning which teams play each other on
each date. No objective function is included; we seek any feasible solution.

\subsection{Sets and Indices}

\begin{align*}
  T &= \text{set of all teams} \quad (|T| = 16) \\
  D &= \text{set of all game dates} \quad (|D| = 16) \\
  i, j &\in T \text{ denote teams } (i \neq j) \\
  d &\in D \text{ denotes a game date}
\end{align*}

\subsection{Parameters}

From Question~1:
\begin{align*}
  H(i) &= \text{set of dates on which team } i \text{ plays at home} \\
  A(i) &= \text{set of dates on which team } i \text{ plays away} \\
  c^H(i,j) &= \text{number of times team } i \text{ hosts team } j \quad \text{[part (b)]} \\
  c^A(i,j) &= \text{number of times team } i \text{ visits team } j\text{'s arena} \quad \text{[part (c)]}
\end{align*}

\subsection{Decision Variables}

Define binary variables:
\[
  x_{ijd} \in \{0,1\} \quad \text{for all } i \neq j \in T,\; d \in D
\]
where $x_{ijd} = 1$ if and only if team $i$ hosts team $j$ on date $d$.

\subsection{Constraints}

\textbf{(e) Home Dates Preserved.}
For each team $i$ and date $d$, team $i$ plays exactly one home game if $d$ is
a home date, and none otherwise:
\[
  \sum_{j \neq i} x_{ijd} = \mathbf{1}[d \in H(i)] \quad \forall\, i \in T,\; d \in D
\]

\textbf{(f) Away Dates Preserved.}
For each team $i$ and date $d$, team $i$ plays exactly one away game if $d$ is
an away date, and none otherwise:
\[
  \sum_{j \neq i} x_{jid} = \mathbf{1}[d \in A(i)] \quad \forall\, i \in T,\; d \in D
\]

\textbf{(g) Home Matchup Counts Preserved.}
Over all dates, team $i$ hosts team $j$ exactly $c^H(i,j)$ times:
\[
  \sum_{d \in D} x_{ijd} = c^H(i,j) \quad \forall\, i \neq j \in T
\]

\textbf{(h) Away Matchup Counts Preserved.}
Over all dates, team $i$ visits team $j$'s arena exactly $c^A(i,j)$ times:
\[
  \sum_{d \in D} x_{jid} = c^A(i,j) \quad \forall\, i \neq j \in T
\]

\subsection{Complete Formulation}

\begin{align*}
  &\textbf{Find} \quad x_{ijd} \in \{0,1\} \quad \forall\, i \neq j \in T,\; d \in D \\[6pt]
  &\textbf{subject to:} \\
  &\quad \sum_{j \neq i} x_{ijd} = \mathbf{1}[d \in H(i)] \quad \forall\, i \in T,\; d \in D \tag{e} \\
  &\quad \sum_{j \neq i} x_{jid} = \mathbf{1}[d \in A(i)] \quad \forall\, i \in T,\; d \in D \tag{f} \\
  &\quad \sum_{d \in D} x_{ijd} = c^H(i,j) \quad \forall\, i \neq j \in T \tag{g} \\
  &\quad \sum_{d \in D} x_{jid} = c^A(i,j) \quad \forall\, i \neq j \in T \tag{h} \\
  &\quad x_{ijd} \in \{0,1\} \quad \forall\, i \neq j \in T,\; d \in D
\end{align*}

\subsection{Model Size and Result}

\begin{itemize}[nosep]
  \item \textbf{Variables:} $16 \times 15 \times 16 = 3{,}840$ binary variables
  \item \textbf{Constraints:} $2 \times 16 \times 16 + 2 \times 16 \times 15 = 992$ linear constraints
\end{itemize}

Gurobi returns \texttt{Status: Optimal}, confirming the constraint system is
consistent and at least one valid reassignment of opponents exists.

% =============================================================================
\section{Question 3 --- Timezone Travel Constraint}
% =============================================================================

\subsection{Additional Constraint}

Building on the feasibility model of Question~2, we add a constraint limiting
how much timezone shift a team can experience across three consecutive games.
For each team $t$ and each triple of consecutive game dates $(d_1, d_2, d_3)$
in its schedule, the sum of consecutive timezone differences must not reach 4 or
more:
\[
  |tz(t,d_2) - tz(t,d_1)| + |tz(t,d_3) - tz(t,d_2)| \leq 3
  \quad \forall\, t \in T,\; \text{consecutive triples } (d_1,d_2,d_3)
\]
where $tz(t,d)$ denotes the UTC offset of the arena where team $t$ plays on
date $d$. Each arena is assigned to one of four timezone groups: Eastern
(UTC$-$5), Central (UTC$-$6), Mountain (UTC$-$7), or Pacific (UTC$-$8).

\subsection{Linearization}

Since $tz(t,d)$ is itself a linear expression in the $x$ variables (each team
plays in exactly one arena per game date), the absolute-value terms are
linearized in the standard way. Let $\delta_{12} = tz(t,d_2) - tz(t,d_1)$ and
$\delta_{23} = tz(t,d_3) - tz(t,d_2)$. Introduce auxiliary variables
$p_{12}, p_{23} \geq 0$ for each (team, triple) pair:
\begin{align*}
  p_{12} &\geq \delta_{12}, \quad p_{12} \geq -\delta_{12} \;\Rightarrow\; p_{12} = |\delta_{12}| \\
  p_{23} &\geq \delta_{23}, \quad p_{23} \geq -\delta_{23} \;\Rightarrow\; p_{23} = |\delta_{23}| \\
  p_{12} &+ p_{23} \leq 3
\end{align*}

\subsection{Result: Infeasibility}

The model is infeasible --- Gurobi proves that no schedule satisfying all
constraints (e)--(h) from Question~2 can simultaneously satisfy the timezone
constraint for every team.

To diagnose this, we compute the IIS (Irreducible Inconsistent Subsystem) ---
the smallest subset of constraints that is still collectively infeasible. The
IIS identifies that the root cause lies with Western Conference teams,
particularly Golden State Warriors, Los Angeles Lakers, and Phoenix Suns, which
are forced by the fixed home/away date structure (constraints e and f) to make
large coast-to-coast timezone jumps within consecutive game triples. Because the
timezone spread from Pacific (UTC$-$8) to Eastern (UTC$-$5) is 3 hours, any
back-to-back crossing of two such time zone boundaries produces a diff sum of 4
or more. The matchup-count constraints (g and h) then prevent any rearrangement
of opponents from resolving all violations simultaneously.

\subsection{Relaxation via Slack Variables}

Since the hard constraint is infeasible, we relax it by introducing a
non-negative slack variable $s(t,k)$ for each team $t$ and consecutive triple
$k$. The constraint becomes:
\[
  p_{12} + p_{23} \leq 3 + s(t,k), \quad s(t,k) \geq 0
\]
We then minimize the total slack across all teams and triples:
\[
  \min \sum_{t,\,k} s(t,k)
\]
This finds the schedule that comes closest to satisfying the timezone constraint
--- the one that violates it by the least possible total amount. The optimal
value is \textbf{27 hours}, distributed across 10 teams and 18 triples. Golden
State Warriors accounts for the largest share (10 of the 27 excess hours across
5 triples), consistent with their position as the westernmost team in the
schedule.

% =============================================================================
\section{Question 4 --- Minimizing Maximum Travel Distance}
% =============================================================================

For Question~4, we considered two different optimization objectives:

\begin{itemize}[nosep]
  \item \textbf{Method 1} minimizes the maximum total travel distance across teams.
  \item \textbf{Method 2} minimizes the maximum deviation from the league-average travel distance.
\end{itemize}

% ── Method 1 ─────────────────────────────────────────────────────────────────
\subsection{Method 1 --- Minimizing Maximum Travel Distance}

\subsubsection*{Objective}

We improve the original schedule by minimizing the maximum total travel distance
among all teams. Introducing a scalar variable $T^*$, the objective is:
\[
  \min\; T^* \quad \text{subject to: } \mathrm{travel}(t) \leq T^* \;\; \forall\, t \in T
\]
This is a minimax formulation: $T^*$ equals the worst-case team travel distance,
and we minimize it. The schedule is otherwise constrained by (e)--(h) from
Question~2.

\subsubsection*{Computing Travel Distance}

Travel distances are computed as great-circle (Haversine) distances in miles
between consecutive arena locations for each team. For team $t$ travelling
between consecutive game dates $d_1$ and $d_2$, four cases arise:

\begin{itemize}[nosep]
  \item \textbf{Home $\to$ Home (HH):} zero travel --- both games at the team's own arena.
  \item \textbf{Home $\to$ Away (HA):} $\mathrm{dist}(\mathrm{arena}(t),\,\mathrm{arena}(h_2))$ where $h_2$ is the host on $d_2$. Linear in $x$.
  \item \textbf{Away $\to$ Home (AH):} $\mathrm{dist}(\mathrm{arena}(h_1),\,\mathrm{arena}(t))$ where $h_1$ is the host on $d_1$. Linear in $x$.
  \item \textbf{Away $\to$ Away (AA):} $\mathrm{dist}(\mathrm{arena}(h_1),\,\mathrm{arena}(h_2))$ --- a product of two binary variables (bilinear).
\end{itemize}

\subsubsection*{Linearization of Away$\to$Away Legs}

The bilinear AA terms are linearized by introducing auxiliary binary variables
$w(t, d_1, d_2, h_1, h_2) \in \{0,1\}$ for each (team, date pair, host pair)
combination:
\[
  w \leq x_{h_1,t,d_1}, \quad w \leq x_{h_2,t,d_2}, \quad w \geq x_{h_1,t,d_1} + x_{h_2,t,d_2} - 1
\]
These three constraints enforce $w = x_{h_1,t,d_1} \cdot x_{h_2,t,d_2}$, the
standard linearization of a binary product. The travel contribution of each AA
leg is then $\mathrm{dist}(\mathrm{arena}(h_1), \mathrm{arena}(h_2)) \cdot w$, which is linear.

\subsubsection*{Model Size}

\begin{itemize}[nosep]
  \item \textbf{Variables:} $3{,}840$ $x$-vars $+$ $\approx 13{,}275$ $w$-vars $+ 1$ $T^* \approx 17{,}116$ binary/continuous variables
  \item \textbf{Constraints:} 992 base $+$ per-team travel bounds $+$ $w$ linearization constraints
\end{itemize}

\subsubsection*{Results}

Gurobi finds an optimal solution with a maximum team travel distance of
\textbf{18,045 miles} (Golden State Warriors), compared to 20,400 miles in the
original schedule --- an \textbf{11.5\% reduction}. The optimized schedule is
saved to \texttt{optimized\_schedule.csv}.

\begin{table}[h]
\centering
\small
\caption{Travel comparison: original vs.\ Method 1 optimized schedule}
\begin{tabular}{lrrr}
\toprule
\textbf{Team} & \textbf{Original (mi)} & \textbf{Optimized (mi)} & \textbf{Savings (mi)} \\
\midrule
Atlanta Hawks         & 7,308  & 7,308  & 0       \\
Boston Celtics        & 15,756 & 17,676 & $-$1,919 \\
Brooklyn Nets         & 12,385 & 14,380 & $-$1,995 \\
Chicago Bulls         & 7,265  & 6,701  & 564     \\
Cleveland Cavaliers   & 9,688  & 13,707 & $-$4,018 \\
Dallas Mavericks      & 13,573 & 13,390 & 183     \\
Denver Nuggets        & 9,922  & 10,406 & $-$484  \\
Golden State Warriors & 20,400 & 18,045 & 2,355   \\
Houston Rockets       & 9,204  & 9,558  & $-$354  \\
Los Angeles Lakers    & 16,166 & 17,712 & $-$1,546 \\
Miami Heat            & 12,798 & 12,486 & 312     \\
Milwaukee Bucks       & 8,458  & 9,060  & $-$603  \\
New York Knicks       & 15,669 & 17,695 & $-$2,026 \\
Philadelphia 76ers    & 13,977 & 14,867 & $-$890  \\
Phoenix Suns          & 16,730 & 11,683 & 5,047   \\
Toronto Raptors       & 10,789 & 11,386 & $-$597  \\
\midrule
\textbf{MAX}          & \textbf{20,400} & \textbf{18,045} & \textbf{2,355} \\
\bottomrule
\end{tabular}
\end{table}

% ── Method 2 ─────────────────────────────────────────────────────────────────
\subsection{Method 2 --- Minimizing Maximum Deviation from League-Average Travel}

\subsubsection*{Objective}

As an alternative improvement criterion, we minimize the maximum deviation of
any team's total travel distance from the league-average travel distance. Let
$\mathrm{travel}(t)$ denote the total travel distance of team $t$, and let
$\bar{D}$ denote the average travel distance across all teams. Introducing a
scalar variable $\delta$, the objective is:
\[
  \min\;\delta
\]
subject to:
\begin{align*}
  \delta &\geq \mathrm{travel}(t) - \bar{D} \quad \forall\, t \in T \\
  \delta &\geq \bar{D} - \mathrm{travel}(t) \quad \forall\, t \in T
\end{align*}
where $\bar{D} = \dfrac{1}{|T|}\displaystyle\sum_{t \in T} \mathrm{travel}(t)$.

This is a min-max fairness formulation: $\delta$ measures the largest absolute
deviation from the league-average travel, and minimizing it makes team travel as
balanced as possible.

\subsubsection*{Model Size}

\begin{itemize}[nosep]
  \item \textbf{Variables:} $3{,}840$ $x$-vars $+$ $\approx 13{,}275$ $w$-vars $+$ 16 travel-vars $+$ 1 $\delta \approx 17{,}132$ variables
  \item \textbf{Constraints:} 992 base $+$ 16 travel-definition $+$ 32 fairness-deviation $+$ $w$ linearization
\end{itemize}

\subsubsection*{Results}

Gurobi finds an optimal solution for Method~2. Under this fairness-based
objective, the optimized schedule achieves a league-average travel distance of
\textbf{13,140 miles} and reduces the maximum deviation from the league average
from 7,895 miles to \textbf{5,000 miles} --- a \textbf{36.7\% improvement in
fairness}. The travel range is also reduced from 13,135 miles to 9,882 miles.

\begin{table}[h]
\centering
\small
\caption{Team-by-team travel comparison: original vs.\ Method 2 optimized schedule}
\begin{tabular}{lrrrr}
\toprule
\textbf{Team} & \textbf{Original} & \textbf{Optimized} & \textbf{vs.\ Orig Avg} & \textbf{vs.\ Opt Avg} \\
\midrule
Golden State Warriors & 20,400 & 18,045 & $+$7,895 & $+$4,905 \\
Phoenix Suns          & 16,730 & 15,532 & $+$4,225 & $+$2,392 \\
Los Angeles Lakers    & 16,166 & 14,909 & $+$3,660 & $+$1,769 \\
Boston Celtics        & 15,756 & 18,140 & $+$3,251 & $+$5,000 \\
New York Knicks       & 15,669 & 17,681 & $+$3,163 & $+$4,542 \\
Philadelphia 76ers    & 13,977 & 14,614 & $+$1,472 & $+$1,474 \\
Dallas Mavericks      & 13,573 & 13,390 & $+$1,067 & $+$251   \\
Miami Heat            & 12,798 & 12,986 & $+$292   & $-$154   \\
Brooklyn Nets         & 12,385 & 13,881 & $-$120   & $+$741   \\
Toronto Raptors       & 10,789 & 11,665 & $-$1,716 & $-$1,474 \\
Denver Nuggets        & 9,922  & 9,376  & $-$2,583 & $-$3,764 \\
Cleveland Cavaliers   & 9,688  & 13,727 & $-$2,817 & $+$587   \\
Houston Rockets       & 9,204  & 11,274 & $-$3,302 & $-$1,866 \\
Milwaukee Bucks       & 8,458  & 8,258  & $-$4,048 & $-$4,881 \\
Atlanta Hawks         & 7,308  & 8,313  & $-$5,198 & $-$4,827 \\
Chicago Bulls         & 7,265  & 8,443  & $-$5,241 & $-$4,697 \\
\midrule
League average        & 12,505 & 13,140 & & \\
\bottomrule
\end{tabular}
\end{table}

\begin{table}[h]
\centering
\small
\caption{Fairness summary: original vs.\ Method 2 optimized schedule}
\begin{tabular}{lrrr}
\toprule
\textbf{Metric} & \textbf{Original} & \textbf{Optimized} & \textbf{Improvement} \\
\midrule
Max deviation from avg ($\delta$) & 7,895  & 5,000  & $+$2,894  \\
Range (max $-$ min travel)        & 13,135 & 9,882  & $+$3,254  \\
Max travel                        & 20,400 & 18,140 & $+$2,260  \\
Min travel                        & 7,265  & 8,258  & $-$994    \\
Total travel                      & 200,087 & 210,232 & $-$10,146 \\
\bottomrule
\end{tabular}
\end{table}

% ── Comparison summary ────────────────────────────────────────────────────────
\subsection{Q4 Comparison Summary}

Method~1 and Method~2 improve the schedule in different ways. Method~1 reduces
the maximum team travel distance, while Method~2 improves fairness by reducing
the gap between teams' travel burdens. In our results, Method~1 performs better
in terms of worst-case travel, whereas Method~2 performs better in terms of
balance across teams. This demonstrates a fundamental trade-off between
efficiency and fairness in schedule design.

\end{document}
